\documentclass[11pt]{article}
\usepackage[utf8]{inputenc}
\usepackage{amsmath,amssymb}
\usepackage[margin=1in]{geometry}
\usepackage{hyperref}
\emergencystretch=3em

\title{State density I: the one-dimensional standard integral at a general scale}
\author{Claude, with Daniel Murfet}
\date{2026-09-06}

\begin{document}
\maketitle
\sloppy

\begin{abstract}
First unit of the state-density programme (the leading coefficient of \texttt{thm:TaylorTree} for
general $\xi,\eta$ as an integral of the fluctuation mass along the exceptional divisor). The
one-dimensional standard integral $I(c)=\int_0^bu^h\eta(u)e^{-\beta(cu^k)^2+\beta cu^k\xi(u)}\,du$ is
studied as a function of an arbitrary scale $c\to\infty$: $c^pI(c)\to\tfrac1kA_{p-1}(\xi(0))\eta(0)$
for data Lipschitz at $0$ and without any sign condition on $\eta(0)$, and
$c^p|I(c)|\le Me^{\beta L^2/2}\tfrac1k\int_0^\infty s^{p-1}e^{-(\beta/2)s^2}ds$ uniformly in $c>0$
when $|\xi|\le L$, $|\eta|\le M$. Zero \texttt{sorry}/\texttt{axiom}.
\end{abstract}

\section{Setting}
Along the divisor the inner one-dimensional integral has scale $c=\sqrt n\prod_jv_j^{k_j}$ depending
on the outer coordinates; the limit theorem is therefore needed in the scale variable, and dominated
convergence in the outer coordinates needs a bound uniform in the scale.

\section{The things formalised}
{\small
\begin{verbatim}
noncomputable def oneDScale (β b c : ℝ) (h k : ℕ) (ξ η : ℝ → ℝ) : ℝ :=
  ∫ u in Ioc 0 b, u ^ h * η u * Real.exp (-β * (c * u ^ k) ^ 2 + β * (c * u ^ k) * ξ u)
theorem oneDScale_tendsto ... :
    Tendsto (fun c => c ^ (((h : ℝ) + 1) / k) * oneDScale β b c h k ξ η) atTop
      (𝓝 (1 / (k : ℝ) * weightedMass β ξ₀ (((h : ℝ) + 1) / k - 1) * η₀))
theorem oneDScale_abs_le ... (hξL : ∀ u ∈ Ioc 0 b, |ξ u| ≤ L)
    (hηM : ∀ u ∈ Ioc 0 b, |η u| ≤ M) :
    c ^ (((h : ℝ) + 1) / k) * |oneDScale β b c h k ξ η|
      ≤ M * Real.exp (β * L ^ 2 / 2) * (1 / (k : ℝ))
        * weightedMass (β / 2) 0 (((h : ℝ) + 1) / k - 1)
\end{verbatim}
}

\section{Proof strategy}
The limit is unit~34's big-$O$ statement (which needs no $\eta(0)\ne0$) turned into a limit of
$n^{p/2}F(n)$ and composed with $n=c^2$; $\sqrt{c^2}=c$ and $(c^2)^{p/2}=c^p$ for $c>0$. The bound
dominates the integrand pointwise by $u^hMe^{\beta L^2/2}e^{-(\beta/2)(cu^k)^2}$ (unit~23's Gaussian
inequality with $t=(cu^k)^2$), then applies the chart power substitution and the scaling lemma of
unit~44 to reduce to the Gaussian moment $\int_0^\infty s^{p-1}e^{-(\beta/2)s^2}$, a
\texttt{weightedMass} at temperature $\beta/2$ and $a=0$.

\section{Roadblocks and resolutions}
\begin{itemize}
\item The Gaussian inequality is stated with $\beta L\sqrt t$; matching $\beta sL$ needed a
\texttt{calc} step with \texttt{ring\_nf} rather than \texttt{rwa}.
\item \texttt{gcongr} on the final product left the inner inequality open; supply
\texttt{mul\_le\_mul\_of\_nonneg\_left} explicitly.
\end{itemize}

\section{Outcome}
Both inputs for integrating the one-dimensional theory along the divisor are in place; the next unit
assembles them by Fubini and dominated convergence for a unique minimal coordinate.
\end{document}
